\documentclass[10pt,twocolumn]{article}

% XeLaTeX. Build: xelatex -> bibtex -> xelatex -> xelatex (see Makefile).
\usepackage{fontspec}
\usepackage{amsmath}
\usepackage[margin=0.85in]{geometry}
\usepackage[numbers,sort&compress]{natbib}
\usepackage{booktabs}
\usepackage{listings}
\usepackage{xcolor}
\usepackage{url}
\usepackage{hyperref}
\hypersetup{colorlinks=true,linkcolor=black,citecolor=black,urlcolor=blue}

\lstset{basicstyle=\ttfamily\footnotesize,columns=fullflexible,breaklines=true,
  frame=none,xleftmargin=0.5em,aboveskip=0.5em,belowskip=0.3em}

% Double-blind toggle for submission.
\newif\ifblind\blindfalse

\title{\bfseries A Universal Positioned-Read Contract for Forensic Evidence:\\
Block-by-Block Filesystem Decoding over Compressed Images,\\
with an Input-Fuzzed Parsing Substrate}

\ifblind
  \author{}
\else
  \author{Albert Hui\\ Security Ronin\\ \texttt{albert.hui@securityronin.com}}
\fi
\date{}

\begin{document}
\maketitle

\begin{abstract}
\noindent
A digital forensic examiner rarely receives a raw disk. Evidence arrives as an
EWF/E01 image, a VMDK or QCOW2 virtual disk, an AFF4 container, a logical AD1 or
DAR archive, or a live device --- and each carries a filesystem (NTFS, ext4,
APFS, HFS+, exFAT, ISO~9660, \dots) that the examiner must walk to find the
artifact that matters. Today every tool re-implements this stack per
container$\times$filesystem pair, and each re-implementation re-derives the same
two error-prone primitives: how to read a byte range from a decoded image, and
how to parse an integer field out of attacker-controllable bytes without
crashing.

We describe a small contract that unifies this surface. Every byte-source ---
however deeply it is compressed, split, or virtualised --- exposes a single
cursorless \emph{positioned read}, \texttt{read\_at(\&self, offset, buf)}; every
filesystem exposes the same inode-addressed tree over that source. We make three
observations that follow from this shape. First, the contract is a genuine
\emph{universal reader}: eighteen fleet formats (six containers, twelve
filesystems and archives) implement one of two traits, and a consumer never
branches on input format. Second, because the read primitive is positioned and
cursorless, walking a scattered on-disk structure (the NTFS \$MFT, whose records
are spread across the volume) over a \emph{block-compressed} container (E01,
stored as independently deflated 32~KiB chunks) decompresses only the chunks that
overlap each small request --- so peak transient decompression memory is
approximately one chunk, independent of image size, rather than the whole
inflated volume. Third, the parsing substrate beneath all of this is a
\texttt{no\_std}, dependency-free, fuzzed \emph{panic-free bounded reader} whose
out-of-range result is $0$, never a crash --- the correct default for untrusted
evidence, and one whose adoption surfaced a real reachable panic (a
denial-of-service) in a shipping carver.

We are careful to separate what is \emph{demonstrated} from what is
\emph{argued}: the end-to-end NTFS-over-E01 composition is validated for
correctness against The Sleuth Kit but exists today as a feature-gated
integration proof, not a shipped pipeline; the memory property is an analytical
bound over a mechanism we verify in code, not yet a measured
resident-set benchmark of an \$MFT walk.
\end{abstract}

\section{Introduction}\label{sec:intro}

Forensic acquisition rarely yields a device the analyst can read directly. The
evidence is an \emph{image} --- a bit-for-bit copy wrapped in a container format
that adds compression, checksums, segmentation, or virtualisation. EWF/E01
\citep{libewf} compresses the media into independently deflated chunks and splits
it across \texttt{.E01, .E02, \dots}\ segments; VMDK and QCOW2 store sparse
grains and copy-on-write clusters; AFF4 \citep{aff4} is a Zip-based map of image
streams; AD1 and DAR are logical file-tree archives, not disks at all. Inside the
decoded media sits a partition table and one or more filesystems, each with its
own on-disk layout.

The classical consequence is a combinatorial explosion of glue: a tool that
mounts NTFS-inside-E01 shares almost no code with one that mounts ext4-inside-VMDK,
even though the only difference is which decoder produces the bytes and which
parser consumes them. Each glue path re-derives two low-level primitives, and
re-derives them \emph{badly}: (i)~how to obtain an arbitrary byte range from the
decoded image, and (ii)~how to read a length or offset field out of
attacker-controllable bytes. The second is a security primitive --- an
out-of-range read of a corrupt image is a crash (a denial of service) or, worse,
a silently wrong value --- and it is copied, subtly differently, into every
reader.

This paper describes a contract that collapses that surface, and two properties
that fall out of its shape.

\subsection*{Contributions}
\begin{enumerate}
  \item \textbf{A universal reader contract} (\S\ref{sec:contract}): two traits
  --- \texttt{ImageSource}, a cursorless positioned-read byte source, and
  \texttt{FileSystem}, an inode-addressed tree over it --- that eighteen distinct
  fleet formats implement, so container decode, partition windowing, and
  filesystem parsing compose without any consumer branching on input format. We
  report the current, \emph{verified} coverage (\S\ref{sec:coverage}), including
  what is \emph{not} yet on the contract.
  \item \textbf{Block-by-block decoding with a one-chunk memory bound}
  (\S\ref{sec:blockbyblock}): because the read primitive is positioned, walking
  the scattered NTFS \$MFT over a block-compressed E01 decompresses only the
  chunks overlapping each $\sim$1~KiB record request into a single reusable
  scratch page --- an analytical bound of $\le K$ chunks for $K$ records, versus
  inflating the entire image. We verify the mechanism against the code and
  validate read correctness against The Sleuth Kit; we are explicit that the
  resident-memory figure is an analytical bound, not yet a benchmark
  (\S\ref{sec:eval}).
  \item \textbf{\texttt{safe-read}, a panic-free-by-construction parsing substrate}
  (\S\ref{sec:saferead}): a \texttt{no\_std}, zero-dependency, fuzzed crate whose
  bounded integer readers return $0$ on any out-of-range offset and never panic.
  Consolidating $\sim$10 hand-rolled copies onto it surfaced a reachable
  out-of-bounds panic --- a denial-of-service on a crafted \$MFT --- in a
  shipping carver, which we fixed under test.
\end{enumerate}

Throughout we label empirical claims by the strength of the check that backs
them, following a two-tier scheme (Table~\ref{tab:tiers}): \emph{Tier~1} is
validated against an independent third-party oracle or real-world data;
\emph{Tier~2} is real engine output whose ground truth is derivable from the
construction or verified in code. We say ``argued'' or ``analytical'' where a
claim is a bound over a verified mechanism rather than a measurement.

\begin{table}[t]\centering\small
\caption{Evidence tiers used in this paper.}\label{tab:tiers}
\begin{tabular}{@{}lp{5.4cm}@{}}
\toprule
Tier & Basis \\
\midrule
1 & Independent third-party oracle (e.g.\ The Sleuth Kit) or real-world data. \\
2 & Real engine output; ground truth derivable from the construction or verified against the source. \\
--- & Analytical bound over a mechanism verified in code, explicitly not measured. \\
\bottomrule
\end{tabular}
\end{table}

\section{The Universal Reader Contract}\label{sec:contract}

The contract has one insight: \emph{everything the examiner opens is a file tree;
the only difference is how deep you decode before files appear.} A bare
filesystem yields files at the root after one parse; a logical archive
(AD1, DAR) is already a file tree; a physical container decodes to a block device
whose partition table becomes directories, each subtree a parsed filesystem.
Two read-only traits express this.

\paragraph{\texttt{ImageSource} --- the byte edge.}
Any decoded media is a read-only, randomly addressable byte stream. The trait
(in \texttt{forensic-vfs}) has one required method:
\begin{lstlisting}
pub trait ImageSource: Send + Sync {
    fn len(&self) -> u64;
    fn read_at(&self, offset: u64,
               buf: &mut [u8]) -> VfsResult<usize>;
    // extents(), view(), source_id(): defaults
}
\end{lstlisting}
There is no cursor and no write method: reads are positioned and parallel-safe by
construction (workers share one \texttt{Arc<dyn ImageSource>}), and evidence is
read-only in the type system --- a mutation is uncompilable. A partition is a
\texttt{SubRange} adapter that adds a base offset and clamps the length, so a read
can never escape its window; container decoders (EWF, VMDK, QCOW2, VHDX, VHD, DMG,
AFF4) present their decoded output as an \texttt{ImageSource}.

\paragraph{\texttt{FileSystem} --- the tree edge.}
A parsed filesystem is an inode-addressed tree, also \texttt{\&self} over interior
mutability so one mounted handle serves $N$ workers. Files are addressed by an
opaque \texttt{FileId} (a filesystem-specific identity, not a raw integer);
directory listings are owned \texttt{Send} streams; content is read by the same
positioned \texttt{read\_at(\&self, id, stream, off, buf)}. Forensically
significant operations --- deleted-inode enumeration, unallocated ranges, named
data streams (NTFS ADS), hard links, slack --- are first-class methods with
default-empty implementations, so a reader implements only what its format
supports.

\paragraph{Composition.}
An engine composes these edges: sniff the input, decode a container to an
\texttt{ImageSource}, detect a volume system and window each partition with a
\texttt{SubRange}, then parse the filesystem in that window to a \texttt{FileSystem}.
A logical archive skips the middle and is a \texttt{FileSystem} over the file
directly. The consumer --- a FUSE mount, a timeliner, a batch carver --- sees one
tree and never learns the input format.

\subsection{Verified coverage, and what is missing}\label{sec:coverage}

Honesty about coverage matters: the contract's value is only as real as the
fraction of formats that actually implement it. Table~\ref{tab:coverage} reports
the status we \emph{verified by enumeration} of the source, not by recollection.
Eighteen formats are on the contract; a substantial set is not yet --- notably
three filesystems (Btrfs, ZFS, ReFS), two container formats (plain VHD, the AFF4
disk side), the entire full-disk-encryption layer (BitLocker, LUKS, FileVault,
VeraCrypt --- a defined \texttt{CryptoLayer} contract with no implementors), and
the temporal/shadow-copy dimension, which the fleet models not as a volume but in
its own \texttt{Snapshot}/\texttt{Store} (raw VSS) and \texttt{TemporalCohort}/
\texttt{EpochTag} (historical-state) vocabulary. We report this gap deliberately:
a ``universal'' reader that silently omitted a third of its formats would be a
worse claim than an honest partial one.

\begin{table}[t]\centering\small
\caption{Contract coverage, verified by source enumeration. ``On contract'' =
implements the named trait today.}\label{tab:coverage}
\begin{tabular}{@{}llp{3.3cm}@{}}
\toprule
Contract & On contract & Not yet \\
\midrule
\texttt{ImageSource} (containers) & ewf, vmdk, qcow2, vhdx, dmg & \textbf{vhd}, \textbf{aff4} (disk) \\
\texttt{FileSystem} (filesystems) & ext4, ntfs, apfs, hfs+, fat/exfat, xfs, udf, iso9660 & \textbf{btrfs}, \textbf{zfs}, \textbf{refs} \\
\texttt{FileSystem} (archives) & zip, dar, ad1 & tar(.gz/.bz2), 7z \\
\texttt{CryptoLayer} (FDE) & --- (none) & bitlocker, luks, filevault, veracrypt \\
Temporal / shadow copies & --- & vsc (\texttt{Snapshot}/\texttt{Store}), \texttt{TemporalCohort}/\texttt{EpochTag} \\
\bottomrule
\end{tabular}
\end{table}

\section{The Parsing Substrate: \texttt{safe-read}}\label{sec:saferead}

Every reader above must extract integers --- offsets, lengths, counts --- from
bytes an adversary may have crafted. The naive idiom, \lstinline|u32::from_le_bytes([b[o],b[o+1],b[o+2],b[o+3]])|,
panics when \texttt{o} runs past the buffer, and \lstinline|o+4| itself can
overflow \texttt{usize}. A panic on a corrupt image is a denial of service; a
tool that instead returns a wrong value is worse. The correct default for
untrusted evidence is: an out-of-range read yields $0$ and never crashes.

We factor this into \texttt{safe-read}, a \texttt{no\_std}, zero-dependency,
\texttt{forbid(unsafe)} crate of six bounded readers
(\texttt{le\_u16/u32/u64}, \texttt{be\_u16/u32/u64}):
\begin{lstlisting}
pub fn le_u32(data: &[u8], off: usize) -> u32 {
    let Some(end) = off.checked_add(4) else { return 0 };
    match data.get(off..end) {
        Some(s) => u32::from_le_bytes(s.try_into()...),
        None => 0,
    }
}
\end{lstlisting}
The guarantee --- ``$0$ out of range, never a panic, for any slice and any
offset including \texttt{usize::MAX}'' --- is exactly a property a fuzz target can
assert, so the crate ships a \texttt{cargo-fuzz} harness that drives arbitrary
(slice, offset) pairs and treats any panic as a failure.

The consolidation was not cosmetic. The fleet contained $\sim$10 hand-rolled
copies of these readers (five in the NTFS USN parser alone, others in the GPT and
live-device layers). Migrating them onto \texttt{safe-read} \emph{surfaced a real
bug}: the NTFS \$MFT carver's local reader indexed directly
(\lstinline|entry[attr_offset+20]|), and a crafted record with
\texttt{attr\_offset} $\in [1005,1015]$ drove that index one byte past the
1024-byte record --- a reachable out-of-bounds panic, i.e.\ a denial-of-service
on a malicious \$MFT. We reproduced it with a failing test, then fixed it by
routing the reads through \texttt{safe-read}: the out-of-range read now returns
$0$, and the record is rejected by the carver's existing content-length guard
rather than crashing the process (Tier~2: a hand-constructed record whose
out-of-range behaviour is derivable; the panic and its removal are verified by
the test).

\section{Block-by-Block Decoding over Compressed Images}\label{sec:blockbyblock}

The contract's second payoff is the paper's central technical claim, and it is a
direct consequence of the read primitive being \emph{positioned} rather than a
cursor.

\paragraph{The problem.}
EWF/E01 stores the acquired media as a sequence of \emph{independently
compressed} chunks --- classically 32~KiB (64 sectors) each, deflated with zlib.
Chunk $i$ covers logical bytes $[\,i\cdot c,\ (i{+}1)\cdot c\,)$ for chunk size
$c$. To read logical byte $x$ one must inflate the chunk $\lfloor x/c\rfloor$; the
compressed chunks are not randomly addressable in the plaintext.

Now consider the NTFS \$MFT: a table of fixed 1024-byte records that describe
every file, but which is itself an ordinary (often fragmented) file, so its
records are scattered across the volume at offsets that must be resolved through
the \$MFT's own data-run list. A timeline or a triage tool reads \emph{many}
records at scattered offsets. The na\"ive way to serve those reads from an E01 ---
decompress the whole image to a temporary file, then seek --- costs
resident/temporary storage equal to the \emph{full inflated media}, which for a
multi-terabyte acquisition is prohibitive.

\paragraph{The mechanism.}
Because \texttt{ImageSource::read\_at} is a bounded byte-range request, the E01
reader can serve it by inflating only the chunks that overlap the range. Its
\texttt{read\_at} loops, and per iteration reads within a single chunk:
\begin{lstlisting}
let chunk_id = (off / self.chunk_size) as usize;
let page = self.read_chunk(chunk_id)?;   // inflate one chunk
let po = (off % self.chunk_size) as usize;
buf[b..b+to_read]
    .copy_from_slice(&page[po..po+to_read]);
\end{lstlisting}
Decompression allocates exactly one chunk-sized scratch page
(\lstinline|vec![0u8; chunk_size]|) and inflates into it, so \emph{peak transient
decompression memory is $\approx$ one chunk, independent of image size}. An NTFS
read of a 1024-byte record therefore resolves its physical offset through the run
list, issues a positioned read, and lands inside a single 32~KiB chunk (32
records share a chunk; a straddle touches two). A partition offset is added by a
\texttt{SubRange} on the way down, transparently.

\paragraph{The bound.}
Reading $K$ \$MFT records touches at most $\sim K$ chunks --- at most
$K\times 32$~KiB inflated \emph{transiently} --- and far fewer in practice, since
32 consecutive records share a chunk and a small LRU chunk cache (the only
interior mutation on the read path) keeps recently touched chunks warm, so
neighbouring records and re-reads pay no repeat decompression. The steady-state
resident footprint is then bounded and \emph{configurable}: \texttt{cache\_size}
chunks (e.g.\ 100 chunks $\approx$ 3.2~MB, 1000 $\approx$ 32~MB), not the inflated
image. The contrast is stark: $O(1)$ transient plus a bounded cache, versus
$O(\text{image})$ to reach records scattered across the disk.

Crucially, this is not an E01-specific hack. The \emph{same} property holds for
any block/cluster/grain-addressed container --- VMDK grains, QCOW2 clusters, VHDX
BAT blocks --- because each implements the identical positioned-read contract.
The optimization is a property of the \emph{contract}, realised once per decoder,
not re-coded per consumer.

\section{Evaluation}\label{sec:eval}

We evaluate three claims at the tiers we can honestly support, and are explicit
about what is not yet measured.

\paragraph{The contracts compose and read correctly (Tier~1).}
An end-to-end integration test opens a real \texttt{.E01}, wraps
\texttt{EwfReader} as an \texttt{ImageSource}, parses NTFS over it, resolves
\texttt{file1.txt} (\$MFT record 37, 408 bytes), and compares the recovered bytes
and metadata against The Sleuth Kit's \texttt{fls}/\texttt{icat} on the same
image. The read path therefore flows small positioned reads down to per-chunk
inflation, and the recovered content matches an independent oracle. \emph{Caveat:}
this composition is gated behind a \texttt{vfs} feature with local path
dev-dependencies --- it is a demonstrated proof, not a shipped product pipeline.
The orchestrating binary (\texttt{disk-forensic}) today uses EWF only for
partition-layout analysis and does not yet mount NTFS over it end-to-end.

\paragraph{The one-chunk memory bound (analytical; mechanism verified).}
We verify in code that decompression allocates a single chunk page and that
\texttt{read\_at} inflates only overlapping chunks; the $\le K\times 32$~KiB
transient bound follows analytically. We do \emph{not} yet report a measured
resident-set benchmark of an \$MFT walk: the repository benchmarks the
\emph{chunk-table} residency (eager vs.\ lazy population --- an eager table is
$\approx\!\text{chunk\_count}\times 16$~B, $\sim$1~GB for a 2~TB image) but not
the peak decompressed RSS of a filesystem walk. We state the memory figure as an
analytical bound over a verified mechanism, and flag a purpose-built measurement
as the natural next step.

\paragraph{Panic-freedom (property, fuzzed) and a real bug (Tier~2).}
\texttt{safe-read}'s never-panic property is asserted by a fuzz target over
arbitrary (slice, offset) inputs. The \$MFT carver panic it surfaced is
reproduced by a deterministic failing test and removed by the fix; the crafted
record's out-of-range behaviour is derivable from the NTFS record layout.

\paragraph{Scope.}
These results establish that the contract \emph{composes correctly} and that the
block-by-block property is \emph{mechanically real}; they do not yet establish an
end-to-end production pipeline nor a measured memory advantage on a large image.
We name both as the outstanding empirical work rather than imply them.

\section{Related Work}\label{sec:related}

Container decoding and filesystem parsing are individually mature: libewf
\citep{libewf} for EWF, The Sleuth Kit \citep{tsk} for NTFS/ext/HFS+, and the
libyal family for many formats. What these provide as separate libraries with
bespoke APIs, our contract unifies behind two traits, so \emph{composition}
(container~$\to$~volume~$\to$~filesystem) is expressed once. dfVFS
\citep{dfvfs} (Plaso) is the closest prior unification --- a Python virtual
file-system layer over many back-ends --- and validates the goal; our contribution
is (i)~a positioned, cursorless, read-only, \texttt{Send+Sync} byte edge that
makes the block-by-block memory property fall out for \emph{any} block-compressed
container, and (ii)~an explicit, fuzzed panic-free parsing substrate as a
first-class dependency rather than an idiom re-copied per reader. The memory
argument is an instance of demand-paged decompression; we make it a property of a
shared contract rather than a per-tool optimization.

\section{Discussion and Limitations}\label{sec:disc}

\paragraph{Threats to validity --- construct.} The ``universal'' claim is bounded
by Table~\ref{tab:coverage}: three filesystems, two containers, all of FDE, and
the shadow-copy dimension are not yet on the contract. We report coverage as
measured, not as aspiration.

\paragraph{Threats to validity --- internal.} The end-to-end composition is a
feature-gated test, and NTFS currently reaches the \texttt{ImageSource} through a
\texttt{Read+Seek} cursor adapter behind a mutex rather than issuing raw parallel
\texttt{read\_at}; this preserves the sliding-window decompression property but
serialises reads and does not yet exploit the contract's inherent concurrency.

\paragraph{Threats to validity --- external.} The memory bound is analytical; a
different container's chunk/cluster size, or a filesystem with larger scattered
reads, changes the constant (``chunks overlapping the range,'' not strictly one).
The LRU cache trades a bounded, configurable resident footprint for avoided
re-decompression --- ``$\approx$ one chunk'' is the \emph{transient} bound, not the
resident one.

\section{Conclusion}\label{sec:concl}

A single positioned-read contract turns a combinatorial glue problem into two
composable edges, and two useful properties follow from its shape rather than from
per-tool effort: scattered filesystem walks over block-compressed images decode
one chunk at a time, and a shared fuzzed parser makes ``don't crash on a corrupt
image'' the default instead of an idiom re-derived --- and mis-derived --- per
reader. We have been deliberate about the boundary between what is demonstrated
(correct composition against an independent oracle; a real panic fixed under test)
and what is argued (the resident-memory advantage, an analytical bound awaiting a
benchmark), and about how much of the format surface is genuinely covered today.
The outstanding work is concrete: finish the format coverage, wire the
composition into a shipping pipeline over a raw positioned-read source, and
measure the \$MFT-walk memory advantage on a multi-terabyte acquisition.

{\small
\bibliographystyle{unsrtnat}
\bibliography{universal-reader}
}

\end{document}
